\section{Introducción}

En un sistema fotovoltaico conectado a la red, la energía que no se produce en un momento dado se toma de la red. En un sistema aislado esa opción no existe: si el sistema se dimensiona mal, las cargas se quedan sin energía. Por eso el banco de baterías es el componente que define la autonomía del sistema. Su capacidad debe cubrir el consumo nocturno y los días de baja radiación, y su tecnología determina cuánto de esa capacidad se puede usar, cuántos ciclos resiste antes de degradarse y cuánto cuesta cada kilovatio-hora que entrega.

El Proyecto~2 consiste en diseñar un sistema fotovoltaico aislado para la casa de Abner López, en Santa Rosa de Oreamuno, Cartago. Para el dimensionamiento se usa un banco de \SI{24}{V} y una profundidad de descarga del 70\,\%, valores fijados en clase como punto de partida. La casa tiene servicio eléctrico, por lo que el sistema se plantea como respaldo de un grupo de cargas críticas.

Este informe describe las tecnologías de baterías que se consideran para el proyecto y los parámetros con que se comparan. Luego identifica opciones comerciales de cada tecnología con su ficha técnica y su precio, las compara en una matriz y selecciona la batería para el proyecto. Finalmente, calcula la capacidad del banco a partir del consumo de las cargas críticas de la casa.
